\section{Analytical solution of the 1D stationary pressure profile}
\label{ann:analytical_pressure}

A one–dimensional steady gas flow through a rigid porous medium is considered.  
The domain is a segment of length $L$, defined for $z\in[0,L]$.  
No volumetric gas sources are present, so the problem is strictly one–dimensional and steady.

The mass flux follows Darcy’s law written for mass transport,
\begin{equation}
\dot{m}'' = \frac{\bar{K}}{\nu}\,\frac{dP}{dz},
\label{eq:darcy}
\end{equation}
where $\bar{K}$ is the intrinsic permeability (m$^2$) and $\nu$ is the kinematic viscosity (m$^2$\,s$^{-1}$).  
The viscosity is assumed inversely proportional to pressure,
\begin{equation}
\nu = \frac{\alpha}{P},
\label{eq:nuP}
\end{equation}
which corresponds to the Chapman--Enskog dilute–gas formulation at constant temperature.

Mass conservation in steady state gives
\begin{equation}
\frac{d\dot{m}''}{dz} = 0,
\end{equation}
implying a spatially uniform mass flux:
\begin{equation}
\dot{m}'' = \text{constant}.
\label{eq:mconst}
\end{equation}

Substituting \eqref{eq:nuP} into Darcy’s law yields
\begin{equation}
\dot{m}'' = \frac{\bar{K} P}{\alpha}\,\frac{dP}{dz},
\label{eq:darcy2}
\end{equation}
so the governing differential equation becomes
\begin{equation}
P\,\frac{dP}{dz} = \frac{\alpha}{\bar{K}}\,\dot{m}''.
\label{eq:ode}
\end{equation}

\subsection{Case 1: fixed pressure at $z=0$ and $z=L$}
\label{ann:analytical_pressure1}
The boundary conditions are
\begin{equation}
P(0)=P_0,
\qquad
P(L)=P_L .
\end{equation}

Integrating \eqref{eq:ode} from $z=0$ to $z=L$ gives
\begin{equation}
\frac{1}{2}\left(P_L^2 - P_0^2\right)
=
\frac{\alpha L}{\bar{K}}\,\dot{m}''.
\end{equation}

The mass flux is therefore
\begin{equation}
\dot{m}'' = \frac{\bar{K}}{2\alpha L}\left(P_L^2 - P_0^2\right).
\label{eq:flux_case1}
\end{equation}

Using \eqref{eq:ode}, the pressure profile becomes
\begin{equation}
P(z)^2
=
P_0^2
+
\left(P_L^2 - P_0^2\right)\frac{z}{L}.
\label{eq:profile_case1}
\end{equation}

\subsection{Case 2: fixed pressure at $z=0$ and imposed mass flux at $z=L$}
\label{ann:analytical_pressure2}
The boundary conditions now read
\begin{equation}
P(0)=P_0,
\qquad
\dot{m}''(L)=\dot{m}''_{\!L},
\end{equation}
where $\dot{m}''_{\!L}$ is the prescribed flux at $z=L$.  
Since the flux is uniform [Eq.~\eqref{eq:mconst}], the entire domain satisfies
\begin{equation}
\dot{m}'' = \dot{m}''_{\!L}.
\label{eq:flux_case2}
\end{equation}

Inserting the imposed flux into \eqref{eq:ode} gives
\begin{equation}
P\,\frac{dP}{dz}
=
\frac{\alpha}{\bar{K}}\dot{m}''_{\!L}.
\end{equation}

Integrating from $0$ to $z$ yields
\begin{equation}
P(z)^2 - P_0^2
=
2\,\frac{\alpha}{\bar{K}}\,\dot{m}''_{\!L}\,z .
\end{equation}

The pressure profile for Case~2 therefore reads
\begin{equation}
P(z)
=
\sqrt{
P_0^2
+
2\,\frac{\alpha}{\bar{K}}\,\dot{m}''_{\!L}\,z
}.
\label{eq:profile_case2}
\end{equation}

If desired, the resulting pressure at $z=L$ follows directly as
\begin{equation}
P(L)
=
\sqrt{
P_0^2
+
2\,\frac{\alpha L}{\bar{K}}\,\dot{m}''_{\!L}
}.
\end{equation}

This concludes the analytical pressure solutions for both a fully prescribed-pressure boundary configuration and a mixed pressure–flux boundary configuration.

